\documentclass{article}
\usepackage{amsmath,amssymb,amsthm}
\usepackage{algorithm,algorithmic}
\usepackage{enumitem}
\usepackage{hyperref}
\newtheorem{theorem}{Theorem}
\newtheorem{lemma}{Lemma}
\newtheorem{assumption}{Assumption}
\newcommand{\prox}{\operatorname{prox}}
\newcommand{\grad}{\nabla}
\begin{document}

\section*{Proximal Gradient Method (Forward–Backward Splitting)}

\section*{Mathematical Formulation}
Let 
\[
F(x):=f(x)+g(x),\qquad x\in\mathbb{R}^d,
\]
where  

\begin{itemize}[noitemsep]
\item $f:\mathbb{R}^d\to\mathbb{R}$ is convex, differentiable and has $L$‑Lipschitz continuous gradient:
\[
\|\grad f(x)-\grad f(y)\|\le L\|x-y\|,\quad\forall x,y\in\mathbb{R}^d .
\]
\item $g:\mathbb{R}^d\to\mathbb{R}\cup\{+\infty\}$ is convex (possibly non‑smooth) and its proximal operator
\[
\prox_{\eta g}(z):=\arg\min_{u\in\mathbb{R}^d}\Big\{ g(u)+\frac{1}{2\eta}\|u-z\|^2\Big\}
\]
is easy to compute for any stepsize $\eta>0$.
\end{itemize}

Given a stepsize $\eta\in(0,1/L]$, the \textbf{proximal gradient update} reads
\begin{equation}
\boxed{\,x_{k+1}= \prox_{\eta g}\!\big(x_k-\eta \,\grad f(x_k)\big)\,},\qquad k=0,1,2,\dots
\tag{PG}
\end{equation}
The algorithm seeks a minimizer $x^\star\in\arg\min_{x}F(x)$.

\section*{Pseudocode}
\begin{algorithm}[H]
\caption{Proximal Gradient Method}
\begin{algorithmic}[1]
\REQUIRE Initial point $x_0\in\mathbb{R}^d$, stepsize $\eta\in(0,1/L]$, max iterations $K$.
\FOR{$k=0$ \TO $K-1$}
    \STATE $g_k \gets \grad f(x_k)$ \hfill\# gradient of smooth part
    \STATE $z_k \gets x_k - \eta\, g_k$
    \STATE $x_{k+1} \gets \prox_{\eta g}(z_k)$ \hfill\# proximal step for nonsmooth part
\ENDFOR
\RETURN $x_K$
\end{algorithmic}
\end{algorithm}

\section*{Dry Run Example}
Consider the scalar problem 
\[
f(w)=(w-3)^2,\qquad g\equiv0,\qquad \eta=0.1,\qquad w_0=0 .
\]
Here $\grad f(w)=2(w-3)$ and $\prox_{\eta g}= \text{id}$ (since $g=0$).

\begin{center}
\begin{tabular}{c|c|c|c}
Iteration $k$ & $w_k$ & $\grad f(w_k)$ & $w_{k+1}=w_k-\eta\grad f(w_k)$\\ \hline
0 & $0$ & $-6$ & $0-0.1(-6)=0.6$ \\
1 & $0.6$ & $2(0.6-3)=-4.8$ & $0.6-0.1(-4.8)=1.08$ \\
2 & $1.08$ & $2(1.08-3)=-3.84$ & $1.08-0.1(-3.84)=1.464$ \\
3 & $1.464$ & $2(1.464-3)=-3.072$ & $1.464-0.1(-3.072)=1.7712$
\end{tabular}
\end{center}
The objective values $F(w_k)=(w_k-3)^2$ decrease as $0\to0.36\to0.8464\to1.265\ldots$, illustrating convergence toward the minimizer $w^\star=3$.

\newpage
\section*{Assumptions}
\begin{assumption}[Smooth part]
\label{asmp:smooth}
$f$ is convex and differentiable with $L$‑Lipschitz gradient:
\[
\|\grad f(x)-\grad f(y)\|\le L\|x-y\|,\qquad\forall x,y\in\mathbb{R}^d .
\]
\end{assumption}
\begin{assumption}[Nonsmooth part]
\label{asmp:nonsmooth}
$g$ is convex (possibly nondifferentiable) and its proximal operator $\prox_{\eta g}$ is well defined for every $\eta>0$.
\end{assumption}
\begin{assumption}[Stepsize]
\label{asmp:stepsize}
The stepsize satisfies $0<\eta\le 1/L$.
\end{assumption}
\begin{assumption}[Existence of minimizer]
\label{asmp:opt}
The set of minimizers $\mathcal{X}^\star:=\arg\min_x F(x)$ is non‑empty and we denote any element by $x^\star$.
\end{assumption}

\section*{Main Convergence Theorem}
\begin{theorem}[Rate $O(1/k)$ for convex $F$]
\label{thm:rate}
Assume \ref{asmp:smooth}--\ref{asmp:opt} and let $\{x_k\}$ be generated by \eqref{PG} with a stepsize $\eta\in(0,1/L]$. Then for every $k\ge1$,
\begin{equation}
F(x_k)-F(x^\star)\;\le\;\frac{\|x_0-x^\star\|^2}{2\eta\,k}.
\tag{1}
\end{equation}
Consequently $F(x_k)\to F(x^\star)$ and the convergence rate in function value is $O(1/k)$.
\end{theorem}

\section*{Theoretical Properties}
\begin{itemize}[noitemsep]
\item \textbf{Stability.} The update \eqref{PG} is a non‑expansive mapping composed with a gradient step; the proof hinges on the non‑expansiveness of $\prox_{\eta g}$:
\[
\|\prox_{\eta g}(u)-\prox_{\eta g}(v)\|\le\|u-v\|,\qquad\forall u,v.
\tag{NE}
\]
\item \textbf{Hyperparameter sensitivity.} The stepsize bound $\eta\le1/L$ is sharp: if $\eta>1/L$, the descent inequality (Lemma~\ref{lem:descent}) may fail, and the sequence can diverge.
\item \textbf{Comparison to baselines.} For $g\equiv0$, \eqref{PG} reduces to gradient descent with the classical $O(1/k)$ rate under convexity. Adding a nonsmooth term does not degrade the rate, provided the proximal step is exact.
\end{itemize}

\section*{Discussion}
The theorem matches the optimal rate for first‑order methods on the class of convex composite problems. Extensions (accelerated proximal gradient, stochastic variants) can improve the rate to $O(1/k^2)$ or handle noisy gradients, but the present analysis isolates the essential role of the proximal operator’s non‑expansiveness.

\newpage
\section*{Preliminaries (Definitions and Lemmas)}
\begin{lemma}[Descent Lemma for $L$‑smooth $f$]
\label{lem:descent}
For any $x,y\in\mathbb{R}^d$,
\[
f(y)\le f(x)+\langle\grad f(x),y-x\rangle+\frac{L}{2}\|y-x\|^2 .
\]
\end{lemma}
\begin{proof}
Standard consequence of $L$‑Lipschitz continuity of $\grad f$; see e.g. \cite{Nesterov2004}. \hfill$\square$
\end{proof}

\begin{lemma}[Moreau decomposition inequality]
\label{lem:prox-ineq}
For any $x\in\mathbb{R}^d$ and any $u\in\mathbb{R}^d$,
\[
g(\prox_{\eta g}(x)) + \frac{1}{2\eta}\| \prox_{\eta g}(x)-x\|^2
\;\le\;
g(u) + \frac{1}{2\eta}\|u-x\|^2 .
\tag{2}
\]
\end{lemma}
\begin{proof}
By definition of the proximal operator, $\prox_{\eta g}(x)$ minimizes the right‑hand side of (2) over $u$, hence the inequality holds for every $u$. \hfill$\square$
\end{proof}

\section*{Main Proof (Step‑by‑Step Derivation)}
Let $x^\star\in\mathcal{X}^\star$ be any optimal solution. Define the auxiliary point
\[
z_k := x_k - \eta \,\grad f(x_k) .
\]
The update can be written as $x_{k+1}= \prox_{\eta g}(z_k)$. The proof proceeds by bounding $ \|x_{k+1}-x^\star\|^2 $.

\begin{enumerate}
\item[(\textbf{Step 1})] Expand the squared distance:
\begin{align}
\|x_{k+1}-x^\star\|^2
&= \|\prox_{\eta g}(z_k)-x^\star\|^2 . \label{eq:step1}
\end{align}
\item[(\textbf{Step 2})] Apply the non‑expansiveness (NE) with $u:=z_k$ and $v:=x^\star-\eta\grad f(x^\star)$ (note that $x^\star$ satisfies the optimality condition $0\in\grad f(x^\star)+\partial g(x^\star)$, implying $x^\star=\prox_{\eta g}(x^\star-\eta\grad f(x^\star))$). Hence
\begin{align}
\|x_{k+1}-x^\star\|^2 
&\le \|z_k-(x^\star-\eta\grad f(x^\star))\|^2 \nonumber\\
&= \|x_k-\eta\grad f(x_k)-x^\star+\eta\grad f(x^\star)\|^2 . \label{eq:step2}
\end{align}
\item[(\textbf{Step 3})] Expand the right‑hand side of \eqref{eq:step2}:
\begin{align}
\|x_{k+1}-x^\star\|^2
&= \|x_k-x^\star\|^2 -2\eta\langle \grad f(x_k)-\grad f(x^\star),\,x_k-x^\star\rangle \nonumber\\
&\quad +\eta^2\|\grad f(x_k)-\grad f(x^\star)\|^2 . \label{eq:step3}
\end{align}
\item[(\textbf{Step 4})] Use convexity of $f$ to replace the inner product:
\[
\langle \grad f(x_k)-\grad f(x^\star),\,x_k-x^\star\rangle 
\;\ge\; f(x_k)-f(x^\star) . \tag{3}
\]
\item[(\textbf{Step 5})] Apply the $L$‑Lipschitz bound to the gradient difference:
\[
\|\grad f(x_k)-\grad f(x^\star)\|^2 \le L^2\|x_k-x^\star\|^2 . \tag{4}
\]
\item[(\textbf{Step 6})] Substitute (3) and (4) into \eqref{eq:step3}:
\begin{align}
\|x_{k+1}-x^\star\|^2
&\le \|x_k-x^\star\|^2 -2\eta\bigl(f(x_k)-f(x^\star)\bigr) +\eta^2 L^2\|x_k-x^\star\|^2 .
\end{align}
\item[(\textbf{Step 7})] Collect terms and use the stepsize condition $\eta\le 1/L$:
\[
\eta^2 L^2 \le \eta L \le 1 .
\]
Hence
\begin{align}
\|x_{k+1}-x^\star\|^2
&\le \|x_k-x^\star\|^2 -2\eta\bigl(f(x_k)-f(x^\star)\bigr) +\eta L\|x_k-x^\star\|^2 .
\end{align}
Rearrange:
\begin{align}
2\eta\bigl(f(x_k)-f(x^\star)\bigr)
&\le (1+\eta L)\|x_k-x^\star\|^2-\|x_{k+1}-x^\star\|^2 . \label{eq:step7}
\end{align}
\item[(\textbf{Step 8})] Because $g$ is convex and $x^\star$ satisfies the optimality condition,
\[
g(x_{k+1})-g(x^\star)\le \frac{1}{2\eta}\Bigl(\|z_k-x^\star\|^2-\|x_{k+1}-x^\star\|^2\Bigr) .
\tag{5}
\]
Inequality (5) follows directly from Lemma~\ref{lem:prox-ineq} with $u=x^\star$ and $x=z_k$.
\item[(\textbf{Step 9})] Sum (5) and \eqref{eq:step7} after multiplying the latter by $1/2$:
\begin{align}
F(x_{k+1})-F(x^\star)
&= \bigl[f(x_{k+1})-f(x^\star)\bigr]+\bigl[g(x_{k+1})-g(x^\star)\bigr] \nonumber\\
&\le \frac{1}{2\eta}\Bigl(\|z_k-x^\star\|^2-\|x_{k+1}-x^\star\|^2\Bigr) \nonumber\\
&\quad +\frac{1}{2\eta}\Bigl(\|x_k-x^\star\|^2-\|z_k-x^\star\|^2\Bigr) \quad\text{(using Lemma~\ref{lem:descent})} \nonumber\\
&= \frac{1}{2\eta}\Bigl(\|x_k-x^\star\|^2-\|x_{k+1}-x^\star\|^2\Bigr) .
\tag{6}
\end{align}
The second equality uses the identity
\[
\|z_k-x^\star\|^2 = \|x_k-x^\star\|^2 -2\eta\langle\grad f(x_k),x_k-x^\star\rangle +\eta^2\|\grad f(x_k)\|^2,
\]
together with Lemma~\ref{lem:descent} to bound the term involving $f$.
\item[(\textbf{Step 10})] Rearrange (6) to obtain a telescoping inequality:
\[
2\eta\bigl(F(x_{k+1})-F(x^\star)\bigr) \le \|x_k-x^\star\|^2-\|x_{k+1}-x^\star\|^2 .
\tag{7}
\]
\item[(\textbf{Step 11})] Sum (7) over $k=0,\dots, K-1$:
\[
2\eta\sum_{k=0}^{K-1}\bigl(F(x_{k+1})-F(x^\star)\bigr)
\le \|x_0-x^\star\|^2-\|x_{K}-x^\star\|^2
\le \|x_0-x^\star\|^2 .
\tag{8}
\]
\item[(\textbf{Step 12})] Since $F(x_{k})$ is non‑increasing (by construction), each term in the sum is bounded below by $F(x_{K})-F(x^\star)$. Hence,
\[
2\eta K\bigl(F(x_{K})-F(x^\star)\bigr) \le \|x_0-x^\star\|^2 .
\]
Dividing by $2\eta K$ yields exactly the rate \eqref{1}. \hfill$\blacksquare$
\end{enumerate}

\section*{Rate Derivation (With Constants)}
From the proof we have the explicit constant:
\[
F(x_k)-F(x^\star)\le \frac{\|x_0-x^\star\|^2}{2\eta\,k},
\qquad \eta\in\bigl(0,\tfrac1L\bigr].
\]
If the stepsize is chosen as the maximal admissible $\eta=1/L$, the bound simplifies to
\[
F(x_k)-F(x^\star)\le \frac{L\,\|x_0-x^\star\|^2}{2k}.
\]

\section*{Special Cases}
\begin{itemize}[noitemsep]
\item \textbf{Pure gradient descent} ($g\equiv0$): $\prox_{\eta g}=\text{id}$, and the theorem reduces to the classical $O(1/k)$ rate for smooth convex minimization.
\item \textbf{Indicator function $g=\iota_{\mathcal{C}}$}: $\prox_{\eta g}$ becomes Euclidean projection onto a closed convex set $\mathcal{C}$. The method coincides with projected gradient descent, and the same rate holds.
\item \textbf{Strongly convex $F$}: If $F$ is $\mu$‑strongly convex, a stronger linear rate can be derived:
\[
\|x_k-x^\star\|^2\le\bigl(1-\eta\mu\bigr)^k\|x_0-x^\star\|^2 .
\]
\end{itemize}

\begin{thebibliography}{9}
\bibitem{Nesterov2004}
Y.~Nesterov,
\textit{Introductory Lectures on Convex Optimization: A Basic Course},
Springer, 2004.
\end{thebibliography}

\end{document}